%========================================================
% INTRODUCTION
%========================================================

As quantum computing technologies continue to advance and quantum communication
links become increasingly relevant for interconnecting remote quantum systems, the
reliable transmission of quantum states between distant nodes is becoming an
increasingly important technological requirement. Optical fibers represent a
natural transmission medium for many such applications because of their
compatibility with existing telecommunications infrastructure and their suitability
for long-distance photonic links
~\cite{wehner_elkouss_hanson_quantum_internet_2018, abadal2024electromagnetic}.

Among the different photonic degrees of freedom, polarization provides a
particularly convenient encoding for quantum information. Polarization-entangled
photon pairs can be generated, manipulated, and measured using established optical
components, making them attractive resources for entanglement distribution and
entanglement-based quantum communication protocols
~\cite{bennett_brassard_mermin_BBM92_1992}. Their transmission through real optical
fibers, however, is affected by birefringence, polarization rotations, phase
accumulation, temporal effects, and environmental fluctuations
~\cite{dowling_2023_polarization_control}. While deterministic local
transformations can remain coherent and preserve entanglement, unresolved
fluctuations or coupling to additional degrees of freedom can produce effective
decoherence and degrade the polarization correlations accessible at the receiver
~\cite{poon_law_stochastic_pmd_2008,
brodsky_loss_polarization_entanglement_2011}.

The preservation of polarization entanglement is particularly relevant to
entanglement-based quantum key distribution. In the BBM92 protocol, an entangled
photon pair is distributed between two distant parties, Alice and Bob, who
independently choose between alternative polarization measurement bases to
establish a secure key
~\cite{bennett_brassard_mermin_BBM92_1992}. Fiber propagation and environmental
perturbations can degrade the distributed correlations, increasing the quantum bit
error rate (QBER) and ultimately limiting the suitability of the link for secure
quantum communication.

BBM92 therefore provides the main operational reference for this thesis, whose
objective is to investigate how source properties and fiber propagation conditions
affect the quality of distributed polarization entanglement and its suitability for
quantum communication. The reference configuration consists of an entangled-photon
source connected to two spatially separated fiber arms, \(A\) and \(B\), with one
photon propagating through each arm and polarization used as the qubit degree of
freedom. The work progressively connects the physical mechanisms affecting this
configuration with their mathematical modeling, numerical simulation, and
operational consequences.

Chapter~\ref{ch:01_mathematical_and_physical_framework} establishes the physical
and mathematical framework required for this analysis. It introduces the
description of polarization qubits and bipartite states, the density-operator
formalism, reduced states, local measurements, and polarization correlations. The
chapter also presents the quantities used throughout the thesis to characterize
state quality, coherence, entanglement, and nonlocality. Particular attention is
devoted to concurrence as the principal entanglement measure adopted in the
numerical analysis. A brief description of the BBM92 protocol and of its main
quantities of interest is also provided, establishing the connection between the
quantum-state description and the operational perspective used in the subsequent
chapters.

Chapter~\ref{ch:02_fiber_channel_modeling} describes how optical fiber propagation
modifies the polarization state through a progressive hierarchy of models with
increasing physical detail. A central distinction is made between coherent
polarization evolution, which can modify measured correlations without necessarily
reducing purity or entanglement, and effective decoherence arising from unresolved
fluctuations or coupling to unobserved degrees of freedom. Reduced models are first
used to isolate coherent phase evolution, statistical dephasing, and local
depolarizing effects in the two propagation arms, providing analytically
interpretable descriptions of the different degradation mechanisms. The analysis
is then extended to a frequency-dependent treatment of polarization-mode
dispersion (PMD)
~\cite{riccardi_polarization_entanglement_distribution_2021,
wai_menyuk_random_birefringence_1996,
antonelli_shtaif_brodsky_pmd_2011}, connecting fiber birefringence and the spectral
properties of the photon pair to the effective degradation of polarization
entanglement. The resulting hierarchy therefore provides a continuous path from
simplified effective channels to a physical description of frequency-dependent
fiber propagation.

Chapter~\ref{ch:03_numerical_experiments_results} implements the developed models
in a modular MATLAB simulation framework and investigates through numerical
experiments how source properties, fiber parameters, PMD, and compensation
conditions affect the distributed entanglement. The analysis aims to identify
regions of parameter space in which a prescribed entanglement quality is
maintained, assess the validity of reduced models against the physical
frequency-dependent description, and relate the resulting entanglement
degradation to BBM92-oriented operational requirements.

Chapter~\ref{ch:04_experimental_activity} complements the theoretical and numerical
work with a preliminary laboratory study of polarization-entanglement
distribution. A simplified two-arm platform based on entangled-photon propagation,
polarization control, single-photon detection, and coincidence measurements is
considered. The objective of this activity is to provide a first laboratory
approach to the measurement conditions and practical requirements relevant to
BBM92-oriented quantum communication over optical fibers. At the present stage,
the experimental activity is therefore intended as an initial exploration of the
physical platform rather than as a complete BBM92 implementation.

Finally, Chapter~\ref{ch:05_conclusions} summarizes the main theoretical,
numerical, and experimental results and discusses the possible extensions of the
developed framework. Particular attention is given to the connection between
physical fiber-channel modeling, the numerical identification of
entanglement-quality regions, and their operational interpretation in the context
of entanglement-based quantum communication.